\documentclass[10pt,a4paper]{amsart}
\usepackage[margin=19mm]{geometry}
\usepackage{amsmath,amssymb,mathtools}
\usepackage[scr=rsfs,cal=euler]{mathalfa}
\usepackage{xfrac}
\usepackage[unicode,hidelinks,bookmarks=false]{hyperref}
\newcommand{\ie}{i.e.\ }
\newcommand{\cf}{cf.\ }
\newcommand{\resp}{respectively\ }
\newcommand{\Subsection}{\subsection}
\providecommand{\nobreakdash}{\nobreak\hbox{-}}
\setlength{\emergencystretch}{2em}
\multlinegap0pt
\usepackage{bm}
\usepackage{bbm}
\usepackage{dsfont}
\usepackage{xspace}
\usepackage{cancel}
\usepackage{mathtools}
\usepackage{amsthm}
\makeatletter
\@ifundefined{lem}{\newtheorem{lem}{Lemma}}{}
\@ifundefined{thm}{\newtheorem{thm}{Theorem}}{}
\makeatother
\newcommand{\Acal}{\mathcal{A}}
\newcommand{\bI}{\mathbf{I}}
\newcommand{\bK}{\mathbf{K}}
\newcommand{\bS}{\mathbf{S}}
\newcommand{\bnu}{\bm{\nu}}
\newcommand{\bu}{\mathbf{u}}
\newcommand{\bvarphi}{\bm{\varphi}}
\newcommand{\bx}{\mathbf{x}}
\newcommand{\td}{\tilde}
\title[Quasi-Minnaert Resonances in 2D Elastic Wave Scattering with]{Quasi-Minnaert Resonances in 2D Elastic Wave Scattering with Applications}
\author{Huaian Diao, Kaixin Lu, Ruixiang Tang, Weisheng Zhou}
\date{Source: arXiv:2505.20768v1 (CC BY 4.0)}
\begin{document}
\maketitle

\noindent\emph{Source and licence.} Quasi-Minnaert Resonances in 2D Elastic Wave Scattering with Applications. Version arXiv:2505.20768v1 (CC BY 4.0), distributed under the Creative Commons Attribution 4.0 International licence, as recorded in the arXiv licence metadata for this exact version and shown on the abstract page https://arxiv.org/abs/2505.20768v1. Full rights notice: licenses/arxiv-2505.20768-CC-BY-4.0.txt.\par\smallskip

\noindent\textit{Editorial scope.} Counted unit: the Thm whose source label is \texttt{thm:isl} in the article (source0.tex lines 525--541, source label \texttt{thm:isl}), delivered together with the article's own complete original proof of it (source0.tex lines 543--864, one \texttt{proof} environment of 322 lines). No mathematics is written, repaired or re-derived: statement and proof are the article's own text line for line. Required context retained uncounted: eq:xtm (lines 212--222); eq:trac (lines 224--226); eq:hicon (lines 250--252); eq:de (lines 258--260); eq:lambda= (lines 267--269); eq:jump (lines 290--292); eq:sol (lines 298--306); eq:or (lines 309--311); eq:hn (lines 390--393); eq:jn (lines 390--393); eq:dtgx (lines 395--399); lem:sinlay (lines 408--446); equ:svt (lines 411--414); lem:qninutw (lines 424--427); eq:qp0 (lines 429--432); lem:qninut (lines 436--439); lem:qni (lines 441--444); lem:npz (lines 447--474); equ:kvt (lines 450--453); eq:ui (lines 477--479); lem:thmxy (lines 483--507); eq:lemk1k2 (lines 485--487); eq:n24 (lines 489--500). Every definition, display and label of those lines is retained unchanged. Omitted from this package and not redistributed: 1--211, 223--223, 227--249, 253--257, 261--266, 270--289, 293--297, 307--308, 312--389, 394--394, 400--407, 415--423, 428--428, 433--435, 440--440, 445--446, 454--476, 480--482, 488--488, 501--524, 542--542, 865--1747. Nothing inside a retained excerpt is omitted or altered: every retained body is delivered exactly as the article's lines are written, so no omission receipt of any kind accompanies this package. Citations: the retained excerpts carry no citation command at all, so no bibliography entry of the article is transcribed and no reference is redistributed. Reference adaptation: none was needed and none was made; every reference inside the delivered unit resolves inside it, so no unresolved reference or citation is produced. Printed slips kept literal: no printed word, symbol, hypothesis or number of the counted statement or of its proof was altered. Layout and class: the article's own class is \texttt{amsart} with packages including graphicx, amssymb, amsthm, epstopdf, geometry, hyperref, booktabs, array, paralist, verbatim, subfig, tabularx, amsmath, amsfonts; the wrapper is the lane's pinned \texttt{amsart} editorial wrapper at 10pt on A4, which restates the theorem-like environments the retained text needs and adds the article's own macro definitions that the retained text needs (\texttt{\textbackslash Acal}, \texttt{\textbackslash bI}, \texttt{\textbackslash bK}, \texttt{\textbackslash bS}, \texttt{\textbackslash bnu}, \texttt{\textbackslash bu}, \texttt{\textbackslash bvarphi}, \texttt{\textbackslash bx}, \texttt{\textbackslash td}), with extra wrapper packages (bm, bbm, dsfont, xspace, cancel, mathtools). The delivered numbering is the unit's own. Assets: no retained span contains a figure environment, a graphics-inclusion command, a PDF-page inclusion, a table or a TikZ picture; no image file is copied into this package, the package declares no asset dependency and no image file is redistributed with it. Licence: version arXiv:2505.20768v1 of this article is distributed under the Creative Commons Attribution 4.0 International licence, as recorded in the arXiv licence metadata for this exact version and shown on the abstract page https://arxiv.org/abs/2505.20768v1. A full rights notice is delivered at licenses/arxiv-2505.20768-CC-BY-4.0.txt. Characters: every retained excerpt is pure ASCII, so the delivered engine, XeLaTeX with its default Latin Modern text font, reports no missing character.
\begin{equation}\label{eq:xtm}
	\left\{
	\begin{array}{ll}
		\mathcal{L}_{\td{\lambda}, \td{\mu}}\bu(\bx) + \omega^2 \td{\rho} \bu(\bx) =0,  &  \bx\in D,  \medskip \\
		\mathcal{L}_{\lambda, \mu}\bu(\bx) + \omega^2\rho\bu(\bx) =0,   &   \bx\in \mathbb{R}^2\backslash \overline{D},  \medskip \\
		\bu(\bx)|_- = \bu(\bx)|_+,     & \bx\in\partial D,  \medskip \\
		\partial_{\td{\bnu}}\bu(\bx)|_- = \partial_{{\bnu}}\bu(\bx)|_+, & \bx\in\partial D,  \medskip \\
		\bu^s:=\bu-\bu^i  \qquad \mbox{satisfies the radiation condition},
	\end{array}
	\right.
\end{equation} 

\begin{equation}\label{eq:trac}
	\partial_{\bnu}\bu=\lambda(\nabla\cdot \bu)\bnu + 2\mu(\nabla^s\bu) \bnu.
\end{equation}

\begin{align}\label{eq:hicon}
	\delta  = \frac{\lambda}{\td{\lambda} }=\frac{\mu}{\td{\mu}}, \quad  \epsilon =\frac{\rho}{\td{\rho}} ,
\end{align}

\begin{align}\label{eq:de}
	\tau= \sqrt{\delta/\epsilon} < 1.
\end{align}

\begin{align}\label{eq:lambda=}
	\lambda=\mathcal{O}(1),\quad\mu=\mathcal{O}(1),\quad\rho=\mathcal{O}(1).
\end{align}    

\begin{align}\label{eq:jump}
	\partial_{\bnu}   \bS_{\partial D}^{\omega}[\bvarphi]|_{\pm}(\bx)=\left( \pm\frac{1}{2}\bI +  \bK_{\partial D}^{\omega, *} \right)[\bvarphi](\bx), \quad \bx\in\partial D,
\end{align}

\begin{align}\label{eq:sol}
	\bu=
	\left\{
	\begin{array}{ll}
		\td{\bS}_{\partial D}^{\omega}[\bvarphi_{1}](\bx), & \bx\in D,  \smallskip \\
		{\bS}_{\partial D}^{\omega}[\bvarphi_{2}](\bx) +\bu^i, &  \bx\in \mathbb{R}^2\backslash \overline{D},
	\end{array}
	\right.
\end{align}

\begin{align}\label{eq:or}
	\Acal(\omega,\delta) [\Phi](\bx)=F(\bx), \quad \bx\in\partial D,
\end{align}

\begin{align}
	J_{n}(t) &=\frac{t^{n} }{2^{n}n! } \left [ 1-\frac{t^{2} }{4(n+1)}  +\frac{t^{4} }{2^{5}(n+2)(n+1) } + \mathcal{O}(t^4)\right ], \quad n\ge 0,\label{eq:jn}\\
	H_{n}(t)&=-\frac{\mathrm{i}2^{n} (n-1)!}{\pi t^{n} } \left [ 1+\frac{t^{2} }{4(n-1)} +\frac{t^{4} }{2^{5}(n-1)(n-2)}+\mathcal{O}(t^4)  \right ], \quad n\ge 4,\label{eq:hn}
\end{align}

\begin{align}
	J_{n}^{\prime}(t)=J_{n-1}(t)-(n/t)J_{n}(t),\quad
	J_{n}^{\prime}(t)=-J_{n+1}(t)+(n/t)J_{n}(t),\label{eq:dtgx}
	%\mathscr{C}_{\nu}^{\prime}(t)&=\frac{1}{2}(\mathscr{C}_{\nu-1}(t)-\mathscr{C}_{\nu+1}(t)),\label{eq:dtgx}
\end{align}

\begin{lem}\label{lem:sinlay}
	%{\normalfont \cite[Theorems 1, Proposition 1 and Lemma 8]{HHZJ}}
	The single-layer potentials $\mathbf{S}_{\partial D_{R}}^{\omega}[e^{\mathrm{i}n\theta}\bm{\upsilon}]$ and $\mathbf{S}_{\partial D_{R}}^{\omega}[e^{\mathrm{i}n\theta}\bm{t}]$ have the following expressions  for $|\bx|=R$:
	\begin{align}
		\mathbf{S}_{\partial D_{R}}^{\omega}\left[e^{\mathrm{i} n \theta}\bm{\upsilon }\right](\bx)&=\alpha_{1 n} e^{\mathrm{i} n \theta} \bm{\upsilon}+\alpha_ {2 n} e^{\mathrm{i} n \theta} \bm{t},\label{equ:svt} \\ 
		\mathbf{S}_{\partial D_{R}}^{\omega}\left[e^{\mathrm{i} n \theta} \bm{t}\right](\bx)&=\alpha_{3 n} e^{\mathrm{i} n \theta} \bm{\upsilon}+\alpha_{4 n} e^{\mathrm{i} n \theta} \mathbf{t},\notag
	\end{align}
	where
	\begin{align}
		\alpha_{1n} &=-\frac{\mathrm{i}\pi}{2\omega^{2}\rho R}\left(n^{2}J_{n}(k_{s}R)H_{n}(k_{s}R)+k_{p}^{2}R^{2}J_{n}^{\prime}(k_{p}R)H_{n}^{\prime}(k_{p}R)\right), \notag\\
		\alpha_{2n} &=\frac{n\pi}{2\omega^{2}\rho}\left(k_{s}J_{n}(k_{s}R)H_{n}^{\prime}(k_{s}R)+k_{p}J_{n}^{\prime}(k_{p}R)H_{n}(k_{p}R)\right), \notag\\
		\alpha_{3n} &=-\frac{n\pi}{2\omega^{2}\rho}\left(k_{s}J_{n}^{\prime}(k_{s}R)H_{n}(k_{s}R)+k_{p}J_{n}(k_{p}R)H_{n}^{\prime}(k_{p}R)\right), \notag\\
		\alpha_{4n} &=-\frac{\mathrm{i}\pi}{2\omega^{2}\rho R}\left(k_{s}^{2}R^{2}J_{n}^{\prime}(k_{s}R)H_{n}^{\prime}(k_{s}R)+n^{2}J_{n}(k_{p}R)H_{n}(k_{p}R)\right). \notag
	\end{align}
	
	The single-layer potentials $\mathbf{S}_{\partial D_{R}}^{\omega}[e^{\mathrm{i}n\theta}\bm{\upsilon}]$ and $\mathbf{S}_{\partial D_{R}}^{\omega}[e^{\mathrm{i}n\theta}\bm{t}]$ have the following expressions for $\bx\in\mathbb{R}^{2}\backslash\overline{D}_{R}$:
	\begin{align}\label{lem:qninutw}
		\mathbf{S}_{\partial D_{R}}^{\omega}[e^{\mathrm{i}n\theta}\bm{\upsilon }](\mathbf{x})&= \frac{-\mathrm{i}\pi}{4\omega^{2}\rho R}\left(nk_{s}RJ_{n}(k_{s}R)\bm{\Psi}_{n}^{s,o}(k_{s}|\mathbf{x}|)+k_{p}^{2}R^{2}J_{n}^{\prime}\big(k_{p}R\big)\bm{\Psi}_n^{p,o}(k_{p}|\mathbf{x}|)\right), \\
		\mathbf{S}_{\partial D_{R}}^{\omega}[e^{\mathrm{i}n\theta}\mathbf{t}](\mathbf{x})&= \frac{-\pi}{4\omega^{2}\rho R}\left(k_{s}^{2}R^{2}J_{n}^{\prime}(k_{s}R)\bm{\Psi}_{n}^{s,o}(k_{s}|\mathbf{x}|)+nk_{p}RJ_{n}(k_{p}R)\bm{\Psi}_n^{p,o}(k_{p}|\mathbf{x}|)\right), \notag
	\end{align}
	where
	\begin{align}
		\bm{\Psi}_{n}^{s,o}(k_{s}|\mathbf{x}|)&= \frac{2nH_{n}(k_{s}|\mathbf{x}|)}{k_{s}|\mathbf{x}|}e^{\mathrm{i}n\theta}\bm{\upsilon }+2\mathrm{i}H_{n}^{\prime}(k_{s}|\mathbf{x}|)e^{\mathrm{i}n\theta}\bm{t},\label{eq:qp0} \\
		\bm{\Psi}_n^{p,o}(k_p|\mathbf{x}|)&= 2H_{n}^{\prime}(k_{p}|\mathbf{x}|)e^{\mathrm{i}n\theta}\bm{\upsilon}+\frac{2\mathrm{i}nH_{n}(k_{p}|\mathbf{x}|)}{k_{p}|\mathbf{x}|}e^{\mathrm{i}n\theta}\bm{t}. \notag
	\end{align}
	Moreover, the function $\bm{\Psi}_{n}^{s,o}(k_{s}|\mathbf{x}|)$  pertains to the s-wave, and  $\bm{\Psi}_n^{p,o}\left(k_p|\mathbf{x}|\right)$ pertains to the p-wave. Both of them are radiating solutions to the equation $(\mathcal{L}_{\lambda,\mu}+\omega^{2}{\rho} )\mathbf{u}=0$ in $\bx\in\mathbb{R}^{2}\backslash\overline{D}_{R}$.
	
	Besides, the single-layer potentials $\mathbf{S}_{\partial D_{R}}^{\omega}[e^{\mathrm{i}n\theta}\bm{\upsilon}]$ and $\mathbf{S}_{\partial D_{R}}^{\omega}[e^{\mathrm{i}n\theta}\bm{t}]$ have the following expressions for $\bx\in B_R$:
	\begin{align}\label{lem:qninut}
		\mathbf{S}_{\partial D_{R}}^{\omega}[e^{\mathrm{i}n\theta}\bm{\upsilon }](\bx)&=\frac{-\mathrm{i}\pi}{4\omega^{2}\rho R}\left(nk_{s}RH_{n}(k_{s}R)\bm{\Psi}_{n}^{s,i}(k_{s}|\bx|)+k_{p}^{2}R^{2}H_{n}^{\prime}\big(k_{p}R\big)\bm{\Psi}_{n}^{p,i}(k_{p}|\mathbf{x}|)\right),\\
		\mathbf{S}_{\partial D_{R}}^{\omega}[e^{\mathrm{i}n\theta}\bm{t}](\bx)&=\frac{-\pi}{4\omega^{2}\rho R}\left(k_{s}^{2}R^{2}H_{n}^{\prime}(k_{s}R)\bm{\Psi}_{n}^{s,i}(k_{s}|\bm{x}|)+nk_{p}RH_{n}(k_{p}R)\bm{\Psi}_{n}^{p,i}(k_{p}|\bm{x}|)\right),\notag
	\end{align}
	where
	\begin{align}
		\bm{\Psi}_{n}^{s,i}(k_{s}|\bm{x}|)&=\frac{2nJ_{n}(k_{s}|\bm{x}|)}{k_{s}|\bm{x}|}e^{\mathrm{i}n\theta}\bm{\upsilon}+2\mathrm{i}J_{n}^{\prime}(k_{s}|\bm{x}|)e^{\mathrm{i}n\theta}\bm{t},\label{lem:qni}\\
		\bm{\Psi}_{n}^{p,i}(k_{p}|\bm{x}|)&=2J_{n}^{\prime}(k_{p}|\bm{x}|)e^{\mathrm{i}n\theta}\bm{\upsilon}+\frac{2\mathrm{i}nJ_{n}(k_{p}|\bm{x}|)}{k_{p}|\bm{x}|}e^{\mathrm{i}n\theta}\bm{t}.\notag
	\end{align}
	Moreover, the function $\bm{\Psi}_{n}^{s,i}(k_{s}|\mathbf{x}|)$  pertains to the s-wave, and the function $\bm{\Psi}_{n}^{p,i}(k_p|\mathbf{x}|)$ pertains to the p-wave. Both of them are entire solutions to the equation $(\mathcal{L}_{\lambda,\mu}+\omega^{2}{\rho} )\mathbf{u}=0$ in $\bx\in{D}_{R}$.
\end{lem}

\begin{lem}\label{lem:npz}
	%{\normalfont \cite[Lemma 9 and Lemma 10]{HHZJ}}
	The N-P operator $\mathbf{K}_{\partial D_{R}}^{\omega,*}$ have the following expressions with two densities $e^{\mathrm{i}n\theta}\bm{\upsilon}$ and $e^{\mathrm{i}n\theta}\bm{t}$,
	\begin{align}
		\mathbf{K}_{\partial D_{R}}^{\omega,*}\left[e^{\mathrm{i}n\theta}\bm{\upsilon}\right]&=a_{1n}e^{\mathrm{i}n\theta}\bm{\upsilon}+a_{2n}e^{\mathrm{i}n\theta}\bm{t},\label{equ:kvt}\\
		\mathbf{K}_{\partial D_{R}}^{\omega,*}\left[e^{\mathrm{i}n\theta}\bm{t}\right]&=b_{1n}e^{\mathrm{i}n\theta}\bm{\upsilon}+b_{2n}e^{\mathrm{i}n\theta}\bm{t},\notag
	\end{align}
	where
	\begin{displaymath}
		a_{1n}=-\frac{1}{2}+g_{1,n}(R),\quad a_{2n}=g_{2,n}(R),\quad b_{1n}=g_{3,n}(R),\quad b_{2n}=-\frac{1}{2}+g_{4,n}(R).
	\end{displaymath}
	The coefficients $g_{i,n}(R)$ ($1 \leq i \leq 4$) can be derived from the tractions $\partial_{\nu}\mathbf{S}_{\partial D_{R}}^{\omega}[e^{\mathrm{i}n\theta}\bm{\upsilon}]\big|_{+}$ and $\partial_{\nu}\mathbf{S}_{\partial D_{R}}^{\omega}[e^{\mathrm{i}n\theta}\bm{t}]\big|_{+}$ evaluated on $\partial D_{R}$, as defined by the following expressions
	\begin{align}
		\partial_{\nu}\mathbf{S}_{\partial D_{R}}^{\omega}[e^{\mathrm{i}n\theta}\bm{\upsilon}]|_{+}&=g_{1,n}(|\bx|)e^{\mathrm{i}n\theta}\bm{\upsilon}+g_{2,n}(|\bx|)e^{\mathrm{i}n\theta}\bm{t},\notag\\
		\partial_{\nu}\mathbf{S}_{\partial D_{R}}^{\omega}[e^{\mathrm{i}n\theta}\bm{t}]|_{+}&=g_{3,n}(|\bx|)e^{in\theta}\bm{\upsilon}+g_{4,n}(|\bx|)e^{in\theta}\bm{t},\notag
	\end{align}
	where
	\begin{align}
		g_{1,n}(|\bx|)=&\frac{\mathrm{i}\pi}{2\omega^{2}\rho R^{2}}\Big(2\mu n^{2}J_{n}(k_{s}R)\big(H_{n}(k_{s}|\bx|)-k_{s}RH_{n}^{\prime}(k_{s}|\bm{x}|)\big)+ \notag\\
		&J_n^{\prime}(k_pR)k_pR(H_n(k_p|\bx|)(\omega^2\rho R^2-2\mu n^2)+2k_p\mu RH_n^{\prime}(k_p|\bx|))\Big), \notag\\
		g_{2,n}(|\bx|)=& - \frac{n\mu\pi}{2\omega^{2}\rho R^{2}}\Big(J_{n}(k_{s}R)H_{n}(k_{s}|\bx|)\Big(k_{s}^{2}R^{2}-2n^{2}\Big)+ \notag\\
		&2R\left(k_{s}J_{n}(k_{s}R)H_{n}^{\prime}(k_{s}|\bx|)+k_{p}J_{n}^{\prime}\left(k_{p}R\right)\left(H_{n}(k_{p}|\bx|)-k_{p}RH_{n}^{\prime}(k_{p}|\bx|)\right)\right)\Big), \notag\\
		g_{3,n}(|\bx|)=& \frac{n\pi}{2\omega^{2}\rho R^{2}}\big(J_{n}\big(k_{p}R\big)H_{n}\big(k_{p}\big|\mathbf{x}\big|\big)\big(\omega^{2}\rho R^{2}-2\mu n^{2}\big)+ \notag\\
		&2\mu R\big(k_pJ_n\big(k_pR\big)H_n'\big(k_p|\bx|\big)+k_sJ_n'(k_sR)\big(H_n(k_s|\bx|)-k_sRH_n'(k_s|\bx|)\big)\big)\big), \notag\\
		g_{4,n}(|\bx|)=& \frac{\mathrm{i}\mu\pi}{2\omega^{2}\rho R^{2}}\big(2n^{2}J_{n}\big(k_{p}R\big)\big(H_{n}\big(k_{p}|\bx|\big)-k_{p}RH_{n}^{\prime}\big(k_{p}|\bm{x}|\big)\big)+ \notag\\
		&J_n^{\prime}(k_sR)k_sR\left(k_s^2R^2H_n(k_s|\bx|)+2k_sRH_n^{\prime}(k_s|\bx|)-2n^2H_n(k_s|\bx|)\right).\notag
	\end{align}
\end{lem}

\begin{equation}\label{eq:ui}
	{\bu}^{i}=\kappa \bm{\Psi}_{n}^{s,i},
\end{equation}

\begin{lem}\label{lem:thmxy}
	Consider the constants $\mathcal{R}_1$ and $\mathcal{R}_2$ as defined in \eqref{eq:R1} and \eqref{eq:R2}. We introduce the constants \(K_1\) and \(K_2\) as:
	\begin{align}\label{eq:lemk1k2}
		K_{1} = \frac{2(1-2\gamma_{1}^2+\gamma_{1}^4)}{\mathcal{R}_{1}}, \quad K_{2} =\frac{2\left(1-2\gamma_{2}^2+\gamma_{2}^4\right)}{\mathcal{R}_{2}},
	\end{align}
	where \(\gamma_1 \in (0, 1)\) and \(\gamma_2 \in (1, R)\) are constants. We formally define the integers \(n_2\) and \(n_4\) as follows:
	\begin{align}\label{eq:n24}
		n_{2}=\left\{
		\begin{array}
			{ll}n_{K_{1}}, &K_{1} > \frac{e^2(\ln\gamma_{1})^2}{4}, \\
			1, & K_{1} \leq \frac{e^2(\ln\gamma_{1})^2}{4},
		\end{array}\right.\quad
		n_{4}=\left\{
		\begin{array}
			{ll}n_{K_{2}}, &K_{2} > \frac{e^2(\ln\gamma_{2})^2}{4}, \\
			1, & K_{2} \leq \frac{e^2(\ln\gamma_{2})^2}{4},
		\end{array}\right.
	\end{align}
	where
	\begin{displaymath}
		n_{K_{1}}= \left\lfloor\frac{2}{\ln\gamma_{1}}W_{-1}\left(\frac{\ln\gamma_{1}}{2\sqrt{K_{1}}}\right)\right\rfloor+1,\quad n_{K_{2}} = \left\lfloor\frac{-2}{\ln\gamma_{2}}W_{-1}\left(\frac{-\ln\gamma_{2}}{2\sqrt{K_{2}}}\right)\right\rfloor+1,
	\end{displaymath}
	with \(W_{-1}\left(\cdot\right)\) denoting a branch of the Lambert \(W\) function.
	Consequently, if \( n \geq n_2 \), then \( K_1 n^2 \leq \gamma_1^{-n} \); similarly, if \( n \geq n_4 \), then \( K_2 n^2 \leq \gamma_2^{n} \).
\end{lem}

\begin{thm}\label{thm:isl}
	Consider the elastic scattering problem described by \eqref{eq:xtm}, which involves an elastic material \(\left(D;\tilde{\lambda},\tilde{\mu},\tilde{\rho}\right)\) embedded in  homogeneous elastic medium \(\left(\mathbb{R}^2 \setminus \overline{D};\lambda,\mu,\rho\right)\), where \(D\) is a unit disk centered at the origin in \(\mathbb{R}^2\). Let \(D_R\) represent a disk of radius \(R\), centered at the origin in \(\mathbb{R}^2\), such that \(D \subseteq D_R\).  Let \(\bu\) denote the internal total wave field in \(D\), and  \(\bu^s\) represent the external scattered field in \(\mathbb{R}^2 \setminus \overline{D}\), both of which satisfy \eqref{eq:xtm} for an incident wave \(\bu^i\) defined in \eqref{eq:ui}. Under the assumptions \eqref{eq:hicon}-\eqref{eq:de} and \eqref{eq:lambda=}, for any \(\gamma_1 \in (0, 1)\), \(\gamma_2 \in (1, R)\), and a sufficiently small \(\varepsilon \in \mathbb{R}_+\), if the index \(n\) associated with \(\bu^i\) satisfies the condition 
	\begin{align}\label{eq:condition}
		n \geq \max(n_1, n_2, n_3, n_4),
	\end{align} where
	\begin{align}\label{eq:n13}
		n_1 = \left\lfloor\frac{\ln\varepsilon}{\ln\gamma_1}\right\rfloor + 1, \quad n_3 = \left\lfloor-\frac{\ln\varepsilon}{\ln\gamma_2}\right\rfloor + 1,
	\end{align}
	with $n_2$ and $n_4$ defined in \eqref{eq:n24}
	then we have
	%\begin{displaymath}
	$$
	\frac{\|\mathbf{u}\|_{L^2(D \setminus \mathcal{N}_{D,1-\gamma_{1}}^{-})^2}^2}{\|\mathbf{u}\|_{L^2(D)^2}^2} \leq \varepsilon + \mathcal{O}\left(\varepsilon\omega^2\right) \ll 1, \quad \frac{\|\mathbf{u}^s\|_{L^2(\left(D_R \setminus \overline{D}\right) \setminus \mathcal{N}_{D,\gamma_{2}-1}^{+})^2}^2}{\|\mathbf{u}^s\|_{L^2(\mathbb{R}^2 \setminus D)^2}^2} \leq \varepsilon + \mathcal{O}\left(\varepsilon\omega^2\right)  { \ll }  1.
	$$
	%\end{displaymath}
	Here, the quantity \(\varepsilon\) characterizes the boundary localization level, while \(\lfloor \cdot \rfloor\) denotes the floor function. 
\end{thm}

\begin{proof}
	We consider the scattering problem described by \eqref{eq:xtm}, with the incident wave defined in \eqref{lem:qni}. According to the expressions for the function \(\bm{\Psi}_{n}^{s,i}\) in \eqref{lem:qni}, one has that on the boundary \(\partial D\):
	\begin{displaymath}
		{\bu}^{i}=f_{1,n}e^{\mathrm{i}n\theta}\bm{\upsilon }+f_{2,n}e^{\mathrm{i}n\theta}\bm{t },
	\end{displaymath}
	where
	\begin{displaymath}
		f_{1,n}=\frac{2n\kappa J_n(k_{s})}{k_{s}},\quad f_{2,n}=2\mathrm{i}\kappa J_{n}'(k_{s}).
	\end{displaymath}
	Moreover, by combining the definition of \(\partial_\nu\) provided in \eqref{eq:trac} with direct computations, we obtain the following on the boundary \(\partial D\):
	\begin{displaymath}
		\partial_\nu\bu^{i}=\tilde{f}_{1,n}e^{\mathrm{i}n\theta}\bm{\upsilon }+\tilde{f}_{2,n}e^{\mathrm{i}n\theta}\bm{t },
	\end{displaymath}
	with \begin{align}
		\tilde{f}_{1,n}&=\kappa \gamma_{1,n},\quad \gamma_{1n}=\frac{4n\mu}{k_s}\left(k_sJ_n'(k_s)-J_n(k_s)\right),\notag\\
		\tilde{f}_{2,n}&=\kappa \gamma_{2,n},\quad \gamma_{2n}=\frac{2\mathrm{i}\mu}{k_s}\left(\left(2n^2-k_s^2\right)J_n(k_s)-2k_sJ_n'(k_s)\right).\notag
	\end{align}
	Hence, the density functions \(\bm{\varphi}_{1}\) and \(\bm{\varphi}_{2}\) under the basis \(\left(e^{\mathrm{i}n\theta}\bm{\upsilon}, e^{\mathrm{i}n\theta}\bm{t}\right)\) can be expressed as
	\begin{align}\label{eq:phizh}
		\bm{\varphi}_1=\varphi_{1,1,n}e^{\mathrm{i}n\theta}\bm{\upsilon }+\varphi_{1,2,n}e^{\mathrm{i}n\theta}\bm{t },\quad\ \bm{\varphi}_2=\varphi_{2,1,n}e^{\mathrm{i}n\theta}\bm{\upsilon }+\varphi_{2,2,n}e^{\mathrm{i}n\theta}\bm{t }.
	\end{align}
	
	Using the jump formula \eqref{eq:jump}, along with \eqref{equ:svt} and \eqref{equ:kvt}, equation \eqref{eq:or} can be reformulated as:
	\begin{align}\label{eq:AXb}\bm{Ax}=\bm{b},\end{align}
	where
	\begin{align}\label{eq:axb}
		\bm{A}=\begin{bmatrix}
			\tilde{\alpha}_{1n} & \tilde{\alpha}_{3n} & -\alpha_{1n} & -\alpha_{3n} \\
			\tilde{\alpha}_{2n} & \tilde{\alpha}_{4n} & -\alpha_{2n} & -\alpha_{4n} \\
			\tilde{g}_{1n}-1 & \tilde{g}_{3n} & -g_{1n} & -g_{3n} \\
			\tilde{g}_{2n} & \tilde{g}_{4n}-1 & -g_{2n} & -g_{4n} \\
		\end{bmatrix},\quad
		\bm{x}=\begin{bmatrix}\varphi _{1,1,n} 
			\\\varphi _{1,2,n} 
			\\\varphi _{2,1,n} 
			\\\varphi _{2,2,n} 		
		\end{bmatrix},
		\quad
		\bm{b}=\begin{bmatrix}f _{1,n} 
			\\f _{2,n}  
			\\\tilde{f}_{1,n} 
			\\\tilde{f}_{2,n}  
		\end{bmatrix}.
	\end{align}
	Here, the parameters $\alpha_{in}$ and $g_{in}$ ($i = 1, 2, 3, 4$) are defined in \eqref{lem:sinlay} and \eqref{lem:npz}, while the parameters $\tilde{\alpha}_{in} $ and $\tilde{g}_{in}$  ($i = 1, 2, 3, 4$) are similarly defined in \eqref{lem:sinlay} and \eqref{lem:npz}, with $(\lambda, \mu)$ replaced by $(\tilde{\lambda}, \tilde{\mu})$.
	
	In the sub-wavelength regime, i.e., $\omega \ll 1$, by combining \eqref{eq:jn}, \eqref{eq:hn}, and \eqref{eq:axb}, the aforementioned equation \eqref{eq:AXb} are equivalently expressed as follows
	
	In the sub-wavelength regime, we perform an asymptotic expansion with respect to \(\omega\), combining the expansions of \(J_n\) and \(H_n\) defined in \eqref{eq:jn} and \eqref{eq:hn}. The aforementioned equation \eqref{eq:AXb} is equivalently expressed as follows:
	\begin{align}\label{eq:xtzh}				\left(\bm{A_0}+\omega^{2}\bm{A_{2}}+\omega^{4}\bm{A_{4}}+\mathcal{O}\left(\frac{\omega^{6}}{n^{6}}\right)\right)\bm{x}=\bm{b},
	\end{align}
	where
	\begin{align}
		\left.\bm{A_0}=\left[\begin{array}{cccc}
			\frac{-n  \tau ^2 (\lambda +3 \mu )}{4 \mu  \left(n^2-1\right) (\lambda +2 \mu )} & \frac{\mathrm{i}  \tau ^2 (\lambda +3 \mu )}{4 \mu  \left(n^2-1\right) (\lambda +2 \mu )} & \frac{n  (\lambda +3 \mu )}{4 \mu  \left(n^2-1\right) (\lambda +2 \mu )} & \frac{-\mathrm{i}  (\lambda +3 \mu )}{4 \mu  \left(n^2-1\right) (\lambda +2 \mu )} \\
			\frac{-\mathrm{i}  \tau ^2 (\lambda +3 \mu )}{4 \mu  \left(n^2-1\right) (\lambda +2 \mu )} & \frac{-n  \tau ^2 (\lambda +3 \mu )}{4 \mu  \left(n^2-1\right) (\lambda +2 \mu )} & \frac{\mathrm{i}  (\lambda +3 \mu )}{4 \mu  \left(n^2-1\right) (\lambda +2 \mu )} & \frac{n  (\lambda +3 \mu )}{4 \mu  \left(n^2-1\right) (\lambda +2 \mu )} \\
			-\frac{1}{2} & -\frac{\mathrm{i}}{2}  \left(1-\frac{\tau ^2 (\lambda +3 \mu )}{\lambda +2 \mu }\right) & -\frac{1}{2} & \frac{-\mathrm{i} \mu }{2 \lambda +4 \mu } \\
			\frac{-\mathrm{i} \mu  \tau ^2}{2 \lambda +4 \mu } & \frac{1}{2} \left(\tau ^2-2\right) & \frac{\mathrm{i} \mu }{2 \lambda +4 \mu } & -\frac{1}{2} \\
		\end{array}\right.\right].\notag
	\end{align}
	\iffalse
	\begin{align}\notag
		\bm{A_{2}}=\begin{bmatrix}
			a_{11} & a_{12} & a_{13}& a_{14} \\
			a_{21} & a_{22} & a_{23}& a_{24} \\
			a_{31} & a_{32} & a_{33}& a_{34}\\
			a_{41} & a_{42} & a_{43}& a_{44}
		\end{bmatrix},\quad\quad
		\bm{A_{4}}=\begin{bmatrix}
			c_{11} & c_{12} & c_{13}& c_{14} \\
			c_{21} & c_{22} & c_{23}& c_{24} \\
			c_{31} & c_{32} & c_{33}& c_{34}\\
			c_{41} & c_{42} & c_{43}& c_{44}
		\end{bmatrix},
	\end{align}
	with	
	\begin{align}	
		a_{11}&=\frac{\rho   \tau ^4 \left(8 \mu ^2+n^2 \left(3 \lambda ^2+12 \lambda  \mu +13 \mu ^2\right)\right)}{-16 \mu ^2 n \left(n^4-5 n^2+4\right) (\lambda +2 \mu )^2},a_{12}=\frac{3 \mathrm{i} \rho  \tau ^4 \left(\lambda ^2+4 \lambda  \mu +5 \mu ^2\right)}{8 \mu ^2 \left(n^4-5 n^2+4\right) (\lambda +2 \mu )^2},\notag\\
		a_{13}&=\frac{\rho   \left(8 \mu ^2+n^2 \left(3 \lambda ^2+12 \lambda  \mu +13 \mu ^2\right)\right)}{16 \mu ^2 n \left(n^4-5 n^2+4\right) (\lambda +2 \mu )^2},\quad a_{14}=\frac{-3 \mathrm{i} \rho   \left(\lambda ^2+4 \lambda  \mu +5 \mu ^2\right)}{8 \mu ^2 \left(n^4-5 n^2+4\right) (\lambda +2 \mu )^2},\notag\\	
		a_{21}&=-\frac{3 \mathrm{i} \rho   \tau ^4 \left(\lambda ^2+4 \lambda  \mu +5 \mu ^2\right)}{8 \mu ^2 \left(n^4-5 n^2+4\right) (\lambda +2 \mu )^2},\quad a_{23}=\frac{3 \mathrm{i} \rho   \left(\lambda ^2+4 \lambda  \mu +5 \mu ^2\right)}{8 \mu ^2 \left(n^4-5 n^2+4\right) (\lambda +2 \mu )^2},\notag\\
		a_{22}&=-\frac{\rho  \tau ^4 \left(\lambda ^2 \left(n^2+8\right)+4 \lambda  \mu  \left(n^2+8\right)+\mu ^2 \left(7 n^2+32\right)\right)}{16 \mu ^2 n \left(n^4-5 n^2+4\right) (\lambda +2 \mu )^2},\notag\\
		a_{24}&=\frac{\rho   \left(\lambda ^2 \left(n^2+8\right)+4 \lambda  \mu  \left(n^2+8\right)+\mu ^2 \left(7 n^2+32\right)\right)}{16 \mu ^2 n \left(n^4-5 n^2+4\right) (\lambda +2 \mu )^2},\notag\\
		a_{31}&=\frac{\rho   \tau ^2 \left(2 \mu  \left(n^2-4\right) (\lambda +2 \mu )-\tau ^2 \left(n^2 \left(3 \lambda ^2+12 \lambda  \mu +17 \mu ^2\right)-8 \mu ^2\right)\right)}{8 \mu  n \left(n^4-5 n^2+4\right) (\lambda +2 \mu )^2},\notag\\
		a_{32}&=\frac{- 2 \mathrm{i} \rho   \tau ^2\mu  \left(n^2-4\right) (\lambda +2 \mu )}{8 \mu  \left(n^4-5 n^2+4\right) (\lambda +2 \mu )^2}\notag\\
		&+\frac{ \mathrm{i} \rho   \tau ^2\tau ^2 \left(\lambda ^2 \left(n^2+2\right)+4 \lambda  \mu  \left(n^2+2\right)+\mu ^2 \left(7 n^2+2\right)\right)}{8 \mu  \left(n^4-5 n^2+4\right) (\lambda +2 \mu )^2},\notag\\
		a_{33}&=\frac{\rho   \left(8 \mu  (\lambda +\mu )+n^2 \left(3 \lambda ^2+10 \lambda  \mu +13 \mu ^2\right)\right)}{8 \mu  n \left(n^4-5 n^2+4\right) (\lambda +2 \mu )^2},\notag\\
		a_{34}&=-\frac{\mathrm{i} \rho   \left(\lambda ^2 \left(n^2+2\right)+2 \lambda  \mu  \left(n^2+8\right)+3 \mu ^2 \left(n^2+6\right)\right)}{8 \mu  \left(n^4-5 n^2+4\right) (\lambda +2 \mu )^2},\notag\\
		a_{41}&=-\frac{\mathrm{i} \left(n^2+2\right) \rho  \tau ^4 \left(\lambda ^2+4 \lambda  \mu +5 \mu ^2\right)}{8 \mu  \left(n^4-5 n^2+4\right) (\lambda +2 \mu )^2},\quad a_{42}=-\frac{3 n \rho   \tau ^4 \left(\lambda ^2+4 \lambda  \mu +5 \mu ^2\right)}{8 \mu  \left(n^4-5 n^2+4\right) (\lambda +2 \mu )^2},\notag\\
		a_{43}&=\frac{\mathrm{i} \left(n^2+2\right) \rho   \left(\lambda ^2+4 \lambda  \mu +5 \mu ^2\right)}{8 \mu  \left(n^4-5 n^2+4\right) (\lambda +2 \mu )^2},\quad a_{44}=\frac{3 n \rho   \left(\lambda ^2+4 \lambda  \mu +5 \mu ^2\right)}{8 \mu  \left(n^4-5 n^2+4\right) (\lambda +2 \mu )^2}.\notag
	\end{align}
	\fi  
	We noticed that 
	$det\left ( A_{0}  \right ) =\frac{ \left(\tau ^2+1\right) (\lambda +3 \mu )^2}{32 \mu ^2 \left(n^2-1\right) (\lambda +2 \mu )^2}\ne0, n\ge 2 $. Through direct calculations, one has 
	\begin{align}\notag
		\bm{A}_{0}^{-1} =\begin{bmatrix}
			b_{11} & b_{12} & \frac{-2}{\tau ^2+1} & \frac{\mathrm{i} \left(1-\tau ^2\right)}{\tau ^2+1} \\
			\frac{2 \mathrm{i} \mu  (\lambda +\mu  (n+2))}{R (\lambda +3 \mu )} & \frac{2 \mu  (\mu +n (\lambda +2 \mu ))}{- (\lambda +3 \mu )} & 0 & -1 \\
			b_{31} & b_{32} & \frac{2}{\tau ^2+1}-2 & \frac{-\mathrm{i} \tau ^2 \left(\tau ^2-1\right)}{\tau ^2+1} \\
			\frac{2 \mathrm{i} \mu  \left(-2 \lambda -4 \mu +\tau ^2 (\lambda +\mu  (n+2))\right)}{ (\lambda +3 \mu )} & \frac{2 \mu  \left(\mu  \tau ^2+n \left(\tau ^2-2\right) (\lambda +2 \mu )\right)}{-(\lambda +3 \mu )} & 0 & -\tau ^2 \\
		\end{bmatrix},
	\end{align}
	with
	\begin{align}
		b_{11}&=\frac{2 \mu  \left(\lambda -\left(\tau ^2 (\lambda +\mu  (n+2))\right)-2 \lambda  n-3 \mu  n\right)}{ \left(\tau ^2+1\right) (\lambda +3 \mu )},\notag\\
		b_{12}&=-\frac{2 \mathrm{i} \mu  \left(2 \lambda +3 \mu +\tau ^2 (\mu +n (\lambda +2 \mu ))-\lambda  n\right)}{ \left(\tau ^2+1\right) (\lambda +3 \mu )},\notag\\
		b_{31}&=\frac{2 \mu  \left(-\left(\tau ^4 (\lambda +\mu  (n+2))\right)+\tau ^2 (\lambda +\mu  n)+2 n (\lambda +2 \mu )\right)}{ \left(\tau ^2+1\right) (\lambda +3 \mu )},\notag\\
		b_{32}&=\frac{2 \mathrm{i} \mu  \left(2 \lambda +4 \mu -\tau ^4 (\mu +n (\lambda +2 \mu ))+\tau ^2 (\mu +\lambda  n)\right)}{ \left(\tau ^2+1\right) (\lambda +3 \mu )}.\notag
	\end{align}
	Thus, the equation \eqref{eq:xtzh} can be reformulated as 
	\begin{align}
		\mathbf{A_0}\left(\mathbf{I}+\omega^2\mathbf{A_0}^{-1}\mathbf{A_2}+\omega^4\mathbf{A_0}^{-1}\mathbf{A_4}+\mathcal{O}\left(\frac{\omega^6}{n^5}\right)\right)\bm{x}=\mathbf{b}.\notag
	\end{align}
	Here, since \(\omega \ll 1\), and the elements of the inverse matrix $\mathbf{A_0}^{-1}$ as well as those of matrix $\mathbf{A_2}$ are bounded, we can obtain that
	\[
	\left\|\omega^2\mathbf{A_0}^{-1}\mathbf{A_2}+\omega^4\mathbf{A_0}^{-1}\mathbf{A_4}+\mathcal{O}\left(\frac{\omega^6}{n^5}\right)\right\|_{\infty}<1.
	\]
	
	
	Consequently, by combining \eqref{eq:jn}, \eqref{eq:hn}, and \eqref{eq:dtgx}, we obtain the asymptotic expansion of the solution as:
	\begin{align}
		\bm{x}=&\left(\mathbf{I}+\omega^2\mathbf{A_0}^{-1}\mathbf{A_2}+\omega^4\mathbf{A_0}^{-1}\mathbf{A_4}+\mathcal{O}\left(\frac{\omega^6}{n^5}\right)\right)^{-1}\mathbf{A_{0}^{-1}b}\notag \\
		%=& \left(\mathbf{I}-\omega^2\mathbf{A_0}^{-1}\mathbf{A_2}+\mathcal{O}\left(\frac{\omega^4}{n^2}\right)\right)\mathbf{\mathbf{A_0}^{-1}\mathbf{b}} \notag\\
		=& \mathbf{A_0^{-1}b}-\omega^2\left(\mathbf{A_0^{-1}A_2 A_0^{-1}b}\right)+\mathcal{O}\left(\frac{\omega^4}{n^2}\left(\mathbf{A_0^{-1}b}\right)\right) \notag\\
		=&\begin{pmatrix}t_{1} 
			\\t_{2} 
			\\t_{3} 
			\\t_{4} 	
		\end{pmatrix} -\omega^2\begin{pmatrix}t_{5}
			\\t_{6}
			\\t_{7}
			\\t_{8}		
		\end{pmatrix}+\mathcal{O}\left(\frac{\omega^{4}}{n^2}\begin{pmatrix}t_{1} 
			\\t_{2} 
			\\t_{3} 
			\\t_{4} 
		\end{pmatrix}\right),\notag
	\end{align}
	where
	\begin{align}
		t_{1}&=2 \kappa \mu  \left(\frac{ 2\left(\tau ^2-3\right) (\lambda +2 \mu )(n-1)\left(  k_s J_n'\left(k_s\right)+ n J_n\left(k_s\right)\right) }{\left(\tau ^2+1\right) (\lambda +3 \mu ) k_s} -\frac{ \left(\tau ^2-1\right)  k_s J_n\left(k_s\right)}{\left(\tau ^2+1\right)}\right),\notag\\
		t_{2}&=-2 \mathrm{i} \kappa\mu  \left(\frac{ 2 (n-1) (\lambda +2 \mu ) (k_s J_n'\left(k_s\right)+nJ_n\left(k_s\right)}{(\lambda +3 \mu ) k_s}- k_s J_n\left(k_s\right)\right),\notag\\
		t_{3}&=4 \kappa\mu  \frac{  (\lambda +2 \mu )  \left((n-1) \tau ^4+(1-3 n) \tau ^2-2\right) k_sJ_n'\left(k_s\right)}{\left(\tau ^2+1\right) (\lambda +3 \mu ) k_s}-2 \kappa\mu  \frac{ \tau ^2 \left(\tau ^2-1\right)  k_s J_n\left(k_s\right)}{\left(\tau ^2+1\right) }\notag\\
		&\quad +4 \kappa\mu  \frac{  n (\lambda +2 \mu ) \left((n-1) \tau ^4-(n-3) \tau ^2+2 n\right) J_n\left(k_s\right)}{\left(\tau ^2+1\right) (\lambda +3 \mu ) k_s},\notag\\
		t_{4}&=-4 \mathrm{i} \kappa \mu  \frac{  (\lambda +2 \mu ) (n-1) \tau ^2(k_sJ_n'\left(k_s\right)+nJ_n\left(k_s\right)) }{(\lambda +3 \mu ) k_s}+2 \mathrm{i} \kappa \mu  \tau ^2  k_s J_n\left(k_s\right)\notag\\
		&\quad-8 \mathrm{i} \kappa \mu  \frac{ n (\lambda +2 \mu ) (k_sJ_n'\left(k_s\right)-J_n\left(k_s\right))}{(\lambda +3 \mu ) k_s}.\notag
	\end{align}
	Accordingly, we derive the coefficients in \eqref{eq:phizh}, one can obtain 
	\begin{align}
		\varphi _{1,1,n}=&\frac{\kappa 2^{3-n} \left(\tau ^2-3\right) (\lambda +2 \mu ) \mu ^{\frac{3}{2}-\frac{n}{2}} \rho ^{\frac{n}{2}-\frac{1}{2}}}{\left(\tau ^2+1\right) (\lambda +3 \mu ) (n-2)!}\omega ^{n-1}+q_{1}\omega ^{n+1}+ \mathcal{O}(\omega^{n+3}),\label{eq:x12} \\
		\varphi _{1,2,n}=&-\frac{\mathrm{i} \kappa 2^{3-n} (\lambda +2 \mu ) \mu ^{\frac{3}{2}-\frac{n}{2}} \rho ^{\frac{n}{2}-\frac{1}{2}}}{(\lambda +3 \mu ) (n-2)!}\omega ^{n-1}+q_{2}\omega ^{n+1}+ \mathcal{O}(\omega^{n+3}),\notag
	\end{align}
	where
	\begin{align}\notag
		q_{1}=\mathcal{O}\left(\frac{\kappa 2^{1-n}  \mu ^{\frac{1}{2}-\frac{n}{2}} \rho ^{\frac{n}{2}+\frac{1}{2}}}{(n-1)!}\right),\quad q_{2}=\mathcal{O}\left(\frac{\mathrm{i} \kappa 2^{1-n}  \mu ^{\frac{1}{2}-\frac{n}{2}} \rho ^{\frac{n}{2}+\frac{1}{2}}}{(n-1)!}\right).
	\end{align}
	
	Based on \eqref{eq:sol}, by substituting the coefficients from \eqref{eq:x12} into formula \eqref{eq:phizh} and combining it with  \eqref{lem:qninut} in Lemma \ref{lem:sinlay}, we obtain
	\begin{align}
		\bu =&\tilde{S}_{\partial D }^{\omega }[\bm{\varphi} _{1} ] (\bx) 
		=\varphi _{1,1,n}\tilde{S}_{\partial D }^{\omega }[e^{\mathrm{i}n\theta }\bm{\upsilon}] (\bx)+\varphi _{1,2,n}\tilde{S}_{\partial D }^{\omega }[e^{\mathrm{i}n\theta }\bm{t}] (\bx) \label{eq:usw}\\
		=&\varphi _{1,1,n}\frac{-\mathrm{i}\pi}{4\omega^{2}\tilde{\rho} }\big(n\tilde{k}_{s}H_{n}(\tilde{k}_{s})\bm{\Psi}_{n}^{s,i}(\tilde{k}_{s}|\bx|)+\tilde{k}_{p}^{2}H_{n}^{\prime}\big(\tilde{k}_{p}\big)\bm{\Psi}_{n}^{p,i}(\tilde{k}_{p}|\mathbf{x}|)\big)\notag\\
		&+\varphi _{1,2,n}\frac{-\pi }{4\omega ^{2} \tilde{\rho}} \left ( n\tilde{k}_{p}H_{n}\left (\tilde{k}_{p} \right ) \bm{\Psi}_{n}^{p,i} \left (\tilde{k}_{p} \left | \bx \right |  \right ) +\tilde{k}_{s}^{2}  H_{n}'\left ( \tilde{k}_{s}  \right )  \bm{\Psi}_{n}^{s,i} \left (\tilde{k}_{s} \left | \bx \right |  \right ) \right ).\notag
	\end{align}
	Subsequently, by substituting  \eqref{lem:qni} from Lemma \ref{lem:sinlay} into \eqref{eq:usw}, applying the asymptotic expansions given in \eqref{eq:jn}, \eqref{eq:hn}, and \eqref{eq:dtgx}, and performing the necessary calculations, one obtains the following asymptotic expansion for $\bu$:
	\begin{align}\label{eq:uae}
		\bu
		%=&\left(\left(\xi _1|\bx|^{n-1}+\xi_2|\bx|^{n+1}\right)e^{\mathrm{i}n\theta }\bm{\upsilon}+\left(\xi _3|\bx|^{n-1}+\xi_4|\bx|^{n+1}\right)e^{\mathrm{i}n\theta }\bm{t}\right)\omega ^{n-1}\\
		%&+\mathcal{O}\left(\xi|\bx|^{n-1}\omega^{n+1}\right)\notag\\
		=&\left(\left(\xi _1|\bx|^{n-1}+\xi_2|\bx|^{n+1}\right)e^{\mathrm{i}n\theta }\bm{\upsilon}+\left(\xi _3|\bx|^{n-1}+\xi_4|\bx|^{n+1}\right)e^{\mathrm{i}n\theta }\bm{t}\right)\omega ^{n-1}\left(1+\mathcal{O}\left(\omega^2\right)\right),
	\end{align}
	where
	\begin{align}
		\xi _1&=\frac{\kappa 2^{1-n} \tau ^2 \mu ^{\frac{1}{2}-\frac{n}{2}} \rho ^{\frac{n}{2}-\frac{1}{2}} \epsilon\left( (\lambda  (n+1)+\mu  (n+5))-\left(n-1\right) \tau ^2 (\lambda +\mu )\right)}{ \left(\tau ^2+1\right) (\lambda +3 \mu ) (n-1)!},\notag\\
		\xi _2&=\frac{\kappa 2^{1-n} \tau ^2 \mu ^{\frac{1}{2}-\frac{n}{2}} \rho ^{\frac{n}{2}-\frac{1}{2}} \epsilon  \left(\tau ^2-1\right)  (n (\lambda +\mu )-2 \mu )}{(n+1) \left(\tau ^2+1\right) (\lambda +3 \mu ) (n-2)!},\notag\\
		\xi _3&=\frac{\mathrm{i}\kappa 2^{1-n} \tau ^2 \mu ^{\frac{1}{2}-\frac{n}{2}} \rho ^{\frac{n}{2}-\frac{1}{2}}\epsilon \left( (\lambda  (n+1)+\mu  (n+5))-\left(n-1\right) \tau ^2 (\lambda +\mu )\right)}{ \left(\tau ^2+1\right) (\lambda +3 \mu ) (n-1)!},\notag\\
		\xi _4&=\frac{\mathrm{i} \kappa 2^{1-n} \tau ^2  \mu ^{\frac{1}{2}-\frac{n}{2}} \rho ^{\frac{n}{2}-\frac{1}{2}} \epsilon (n-1)\left(\tau ^2-1\right)(\lambda  (n+2)+\mu  (n+4))}{(n+1) \left(\tau ^2+1\right) (\lambda +3 \mu )  (n-1)!}.\notag
	\end{align}
	Furthermore, through rigorous computations, one can obtain
	\begin{align}
		\xi _1&+\xi _2+\xi _3+\xi _4=\frac{(1+\mathrm{i}) \kappa 2^{1-n} \tau ^2 \epsilon  \mu ^{\frac{1-n}{2}} \rho ^{\frac{n-1}{2}} \left(\mathrm{i} (n-1) \tau ^2+(2-\mathrm{i}) n+(2+\mathrm{i})\right)}{(n+1) \left(\tau ^2+1\right)  (n-1)!}.\notag
		%\xi&=\frac{\kappa2^{-n-2} \tau ^2 \epsilon  \mu ^{-\frac{n}{2}-\frac{1}{2}} \rho ^{\frac{n}{2}+\frac{1}{2}}}{(n-1)!}.\notag 
	\end{align}		
	Based on \eqref{eq:uae}, we can obtain 
	\begin{align}
		\|\bu\|_{L^2(D \setminus \mathcal{N}_{D,1-\gamma_{1}}^{-})^2}^2
		%=&\int_{D \setminus \mathcal{N}_{D,1-\gamma_{1}}^{-}} \left(\left ( \xi_{1}^{2} - \xi_{3}^{2}\right )  | \bx | ^{2n-2}+\left ( 2\xi_{1}\xi_{2}- 2\xi_{3}\xi_{4}\right )|\bx| ^{2n} +\left ( \xi_{2}^{2} - \xi_{4}^{2}\right )|\bx|^{2n+2}\right)\notag\\
		%&\times\omega^{2n-2}\left(1+ \mathcal{O}(\omega^{2})\right)\mathrm{d}\bx\notag\\
		%=&\int_0^{2\pi}\int_0^{\gamma_1} \Big\{\left (\left ( \xi_{1}^{2} - \xi_{3}^{2}\right ) r ^{2n-2}+\left ( 2\xi_{1}\xi_{2}- 2\xi_{3}\xi_{4}\right )r^{2n} +\left ( \xi_{2}^{2} - \xi_{4}^{2}\right ) r ^{2n+2}\right )\omega^{2n-2}\notag\\
		%&\times\left(1+ \mathcal{O}(\omega^{2})\right)\Big\}r\mathrm{d}r\mathrm{d}\theta\notag\\
		% =& \left(\frac{2\pi (\xi_{1}^{2}- \xi_{3}^{2})}{2n} \gamma_1 ^{2n}+\frac{2\pi (2\xi_{1}\xi_{2}- 2\xi_{3}\xi_{4})}{2n+2} \gamma_1 ^{2n+2}+\frac{2\pi (\xi_{2}^{2}- \xi_{4}^{2})}{2n+4} \gamma_1 ^{2n+4}\right)\notag\\
		% &\times\omega^{2n-2}\left(1+ \mathcal{O}(\omega^{2})\right)\notag\\
		=&\left(\zeta _{1,n} \gamma_1 ^{2n}+\zeta _{2,n} \gamma_1 ^{2n+2}+\zeta _{3,n} \gamma_1 ^{2n+4}\right)\omega ^{2n-2}\left(1+\mathcal{O}(\omega^{2})\right)\notag\\
		=&\zeta _{1,n}\gamma_1 ^{2n}\left(1+\zeta _{4,n} \gamma_1 ^{2}+\zeta _{5,n} \gamma_1 ^{4}\right)\omega ^{2n-2}\left(1+\mathcal{O}(\omega^{2})\right),\label{eq:unorm}
	\end{align}
	where
	\begin{align}
		&\zeta _{1,n}=\frac{\pi  \kappa^2 2^{3-2 n} \tau ^4 \mu ^{1-n} \rho ^{n-1} \epsilon^{2}\left(\lambda +5 \mu -\left((n-1) \tau ^2 (\lambda +\mu )\right)+\lambda  n+\mu  n\right)^2}{n \left(\tau ^2+1\right)^2 (\lambda +3 \mu )^2  (n-1)!^2},\label{eq:zetau}\\
		&\zeta_{2,n}=\frac{\left(\lambda +5 \mu -(n-1) \tau ^2 (\lambda +\mu )+\lambda  n+\mu  n\right)}{(n+1) \left(\tau ^2+1\right)^2 (\lambda +3 \mu )^2(n-2)!(n-1)!}\notag\\
		&\quad\quad\times \pi  \kappa^2 4^{2-n} \tau ^4 \left(\tau ^2-1\right) \epsilon ^2 (\lambda +\mu ) \mu ^{1-n} \rho ^{n-1},\notag\\
		&\zeta_{3,n}=\frac{\pi  \kappa^2 4^{1-n} (n-1)^2 \tau ^4 \left(\tau ^2-1\right)^2 \epsilon ^2 \mu ^{1-n} \rho ^{n-1} }{(n+1)^2 (n+2) \left(\tau ^2+1\right)^2 (\lambda +3 \mu )^2  (n-1)!^2}((\lambda  (n+2)+\mu  (n+4))^2\notag\\
		&\quad\quad+(n (\lambda +\mu )-2 \mu )^2),\notag\\
		&\zeta _{4,n}=\frac{-2 (n-1) n \left(\tau ^2-1\right) (\lambda +\mu )}{\left(n^2-1\right) \tau ^2 (\lambda +\mu )-(n+1) (\lambda  (n+1)+\mu  (n+5))}=-2+\mathcal{O}\left(\frac{1}{n}\right),\notag\\
		&\zeta _{5,n}=\frac{\left(\lambda ^2 (n^2+2n+2)+2 \lambda  \mu  (n (n+2)+4)+\mu ^2 (n^2+2n+10)\right)}{(n+2) \left((n+1) (\lambda  (n+1)+\mu  (n+5))-\left(n^2-1\right) \tau ^2 (\lambda +\mu )\right)^2}\notag\\
		&\quad\quad\times (n-1)^2 n \left(\tau ^2-1\right)^2=1+\mathcal{O}\left(\frac{1}{n}\right),\notag
	\end{align}
	By substituting the values of \(\zeta _{4,n}\) and \(\zeta _{5,n}\), we can derive	
	\begin{align}\label{eq:zetar}
		1+\zeta _{4,n}+\zeta _{5,n}=\frac{\mathcal{R}_{1}}{n^2}+\mathcal{O}\left(\frac{1}{n^3}\right),
	\end{align}	
	where
	\begin{align}\label{eq:R1}
		\mathcal{R}_{1}=\frac{\lambda ^2 \left(3 \tau ^4-10 \tau ^2+11\right)+2 \lambda  \mu  \left(5 \tau ^4-18 \tau ^2+25\right)+\mu ^2 \left(11 \tau ^4-34 \tau ^2+59\right)}{\left(\tau ^2-1\right)^2 (\lambda +\mu )^2}.
	\end{align}
	Consequently, based on \eqref{eq:zetau} and \eqref{eq:zetar}, one can obtain 
	\begin{align}
		\frac{1+\zeta _{4,n} \gamma_1 ^{2}+\zeta _{5,n} \gamma_1 ^{4}}{1+\zeta _{4,n}+\zeta _{5,n}}&=\frac{1+\left(-2+\mathcal{O}\left(\frac{1}{n}\right)\right)\gamma_{1}^{2}+\left(1+\mathcal{O}\left(\frac{1}{n}\right)\right)\gamma_{1}^{4}}{\frac{\mathcal{R}_{1}}{n^2}+\mathcal{O}(\frac{1}{n^3})}\notag\\
		&=\frac{n^2}{\mathcal{R}_{1}}\left(1-2\gamma_{1}^2+\gamma_{1}^4+\mathcal{O}\left(\frac{1}{n}\right)\right)\cdot\left(1+\mathcal{O}\left(\frac{1}{n}\right)\right)\notag\\
		&=\frac{1-2\gamma_{1}^2+\gamma_{1}^4}{\mathcal{R}_{1}} n^2+\mathcal{O}\left(n\right)<2 \times \frac{1-2\gamma_{1}^2+\gamma_{1}^4}{\mathcal{R}_{1}} n^2,\notag
	\end{align}  
	where $K_{1} = \frac{2\left(1 - 2\gamma_{1}^2 + \gamma_{1}^4\right)}{\mathcal{R}_{1}}$ is defined in \eqref{eq:lemk1k2}. According to Lemma \ref{lem:thmxy}, we obtain
	\[\frac{1+\zeta _{4,n} \gamma_1 ^{2}+\zeta _{5,n} \gamma_1 ^{4}}{1+\zeta _{4,n}+\zeta _{5,n}}<K_{1}n^2\leq\gamma_1^{-n}.\]
	Consequently, the energy distribution of the total internal field on the boundary can be evaluated, and direct calculations show that the result holds when $n$ satisfies \eqref{eq:condition},
	\begin{align}
		\frac{\|\mathbf{u}\|_{L^2(D \setminus \mathcal{N}_{D,1-\gamma_{1}}^{-})^2}^2}{\|\mathbf{u}\|_{L^2(D)^2}^2}=
		&\frac{\zeta _{1,n}(1 +\zeta _{4,n} \gamma_1 ^{2}+\zeta _{5,n} \gamma_1 ^{4})\omega ^{2n-2}\gamma_1 ^{2n}\left(1+\mathcal{O}( \omega^{2})\right)}{\zeta _{1,n}(1 +\zeta _{4,n} +\zeta _{5,n} )\omega ^{2n-2}\left(1+\mathcal{O}(\omega^{2})\right)}\notag\\
		\leq&\gamma_1^{n}\left(1+\mathcal{O}\left(\omega^2\right)\right)\left(1-\mathcal{O}\left(\omega^2\right)\right)\leq\varepsilon+\mathcal{O}\left(\varepsilon\omega^2\right)\ll1.\notag
	\end{align}
	Then, we can conclude that the energy of the internal scattering field $\bu$ in the elastic scattering problem \eqref{eq:xtm} is predominantly localized on the inner surface of $D$.
	With a calculation similar to \eqref{eq:x12}, one can obtain the following result for the coefficients in \eqref{eq:phizh}
	\begin{align}\label{eq:phi212}
		\varphi _{2,1,n}=&\frac{\kappa2^{3-n} \left(\tau ^2-1\right)^2 (\lambda +2 \mu ) \mu ^{\frac{3}{2}-\frac{n}{2}} \rho ^{\frac{n}{2}-\frac{1}{2}}}{\left(\tau ^2+1\right) (\lambda +3 \mu ) (n-2)!}\omega ^{n-1}+q_{3}\omega ^{n+1}+ \mathcal{O}(\omega^{n+3}),\\
		\varphi _{2,2,n}=&-\frac{\mathrm{i} \kappa 2^{3-n} \left(\tau ^2-1\right) (\lambda +2 \mu ) \mu ^{\frac{3}{2}-\frac{n}{2}} \rho ^{\frac{n}{2}-\frac{1}{2}}}{(\lambda +3 \mu ) (n-2)!}\omega ^{n-1}+q_{4}\omega ^{n+1}+ \mathcal{O}(\omega^{n+3}), \notag
	\end{align}
	where
	\begin{align}
		q_{3}=\mathcal{O}\left(\frac{\kappa 2^{1-n}  \mu ^{\frac{1}{2}-\frac{n}{2}} \rho ^{\frac{n}{2}+\frac{1}{2}}}{(n-1)!}\right),\quad q_{4}=\mathcal{O}\left(\frac{\mathrm{i} \kappa 2^{1-n}  \mu ^{\frac{1}{2}-\frac{n}{2}} \rho ^{\frac{n}{2}+\frac{1}{2}}}{(n-1)!}\right).\notag
	\end{align}
	Based on \eqref{eq:sol}, by substituting the coefficients from \eqref{eq:phi212} into formula \eqref{eq:phizh} and incorporating  \eqref{lem:qninutw} from Lemma \ref{lem:sinlay}, we obtain 
	\begin{align}
		\bu^s =&S_{\partial D }^{\omega }[\bm{\varphi} _{2} ] (\bx) 
		=\varphi _{2,1,n}S_{\partial D }^{\omega }[e^{\mathrm{i}n\theta }\bm{\upsilon}] (\bx)+\varphi _{2,2,n}S_{\partial D }^{\omega }[e^{\mathrm{i}n\theta }\mathbf{t}] (\bx) \label{eq:us}\\
		=&\varphi _{2,1,n} \frac{-\mathrm{i}\pi}{4\omega^{2}\rho }\big(nk_{s}J_{n}(k_{s})\bm{\Psi}_{n}^{s,o}(k_{s}|\mathbf{x}|)+k_{p}^{2}J_{n}^{\prime}\big(k_{p}\big)\bm{\Psi}_{n}^{p,o}(k_{p}|\mathbf{x}|)\big) \notag\\
		&+\varphi _{2,2,n} \frac{-\pi}{4\omega^{2}\rho }\big(k_{s}^{2}J_{n}^{\prime}(k_{s})\bm{\Psi}_{n}^{s,o}(k_{s}|\mathbf{x}|)+nk_{p}J_{n}(k_{p})\bm{\Psi}_{n}^{p,o}(k_{p}|\mathbf{x}|)\big).\notag
	\end{align}
	Subsequently, by substituting \eqref{eq:qp0} from Lemma \ref{lem:sinlay} into \eqref{eq:us}, applying the asymptotic expansions in \eqref{eq:jn}, \eqref{eq:hn}, and \eqref{eq:dtgx}, and carrying out the necessary calculations, one can obtain that $\bu^s$ have the following asymptotic expansion: 
	\begin{align}
		\bu^{s}
		%=&\left(\left(\xi _5|\bx|^{-n-1}+\xi_6|\bx|^{-n+1}\right)e^{\mathrm{i}n\theta }\bm{\upsilon}+\left(\xi _7|\bx|^{-n-1}+\xi_8|\bx|^{-n+1}\right)e^{\mathrm{i}n\theta }\bm{t}\right)\omega ^{n-1}\notag\\
		%&+\mathcal{O}\left(\xi^s|\bx|^{-n-1}\omega^{n+1}\right)\notag\\
		=&\left(\left(\xi _5|\bx|^{-n-1}+\xi_6|\bx|^{-n+1}\right)e^{\mathrm{i}n\theta }\bm{\upsilon}+\left(\xi _7|\bx|^{-n-1}+\xi_8|\bx|^{-n+1}\right)e^{\mathrm{i}n\theta }\bm{t}\right)\omega ^{n-1}\notag\\
		&\times \left(1+\mathcal{O}\left(\omega^2\right)\right),\notag
	\end{align}
	where
	\begin{align}
		\xi _5&=-\frac{\kappa 2^{1-n} \mu ^{\frac{1}{2}-\frac{n}{2}} \rho ^{\frac{n}{2}-\frac{1}{2}} (n-1) \left(\tau ^2-1\right)  \left(\tau ^2 (\lambda +3 \mu )+(n+1) (\lambda +\mu )\right)}{(n+1) \left(\tau ^2+1\right) (\lambda +3 \mu )  (n-1)!},\notag\\
		\xi _6&=\frac{\kappa 2^{1-n} \mu ^{\frac{1}{2}-\frac{n}{2}} \rho ^{\frac{n}{2}-\frac{1}{2}} \left(\tau ^2-1\right) (2 \mu +n (\lambda +\mu ))}{\left(\tau ^2+1\right) (\lambda +3 \mu )  (n-1)!},\notag\\
		\xi _7&=\frac{\mathrm{i} \kappa 2^{1-n}\mu ^{\frac{1}{2}-\frac{n}{2}} \rho ^{\frac{n}{2}-\frac{1}{2}} (n-1) \left(\tau ^2-1\right)  \left(\tau ^2 (\lambda +3 \mu )+(n+1) (\lambda +\mu )\right)}{(n+1) \left(\tau ^2+1\right) (\lambda +3 \mu ) (n-1)!},\notag\\
		\xi _8&=-\frac{\mathrm{i}  \kappa 2^{1-n}\mu ^{\frac{1}{2}-\frac{n}{2}} \rho ^{\frac{n}{2}-\frac{1}{2}} \left(\tau ^2-1\right)  (-2 (\lambda +2 \mu )+n (\lambda +\mu ))}{\left(\tau ^2+1\right) (\lambda +3 \mu ) (n-1)!}.\notag
	\end{align}
	Additionally, through rigorous computations, one has
	\begin{align}
		\xi _5&+\xi _6+\xi _7+\xi _8=\frac{(1+\mathrm{i} ) \kappa 2^{1-n} \left(\tau ^2-1\right) \mu ^{\frac{1}{2}-\frac{n}{2}} \rho ^{\frac{n}{2}-\frac{1}{2}} \left(\mathrm{i}  (n-1) \tau ^2+n+1\right)}{(n+1) \left(\tau ^2+1\right)(n-1)!}.\notag
		%\xi ^{s}&=\mathcal{O}\left(\frac{\kappa 2^{-n-2} \mu ^{-\frac{n}{2}-\frac{1}{2}} \rho ^{\frac{n}{2}+\frac{1}{2}} }{(n-1)!}\right).\notag
	\end{align}
	Similar to \eqref{eq:unorm}, we can obtain
	\begin{align}
		&\left\|\bu^{s}\right\|_{L^{2}\left(\left(D_R \setminus \overline{D}\right) \setminus \mathcal{N}_{D,\gamma_{2}-1}^{+}\right)^{2}}^{2}\notag\\
		%=&\int_{\left(D_R \setminus \overline{D}\right) \setminus \mathcal{N}_{D,\gamma_{2}-1}^{+}} \left(\left ( \xi_{5}^{2} - \xi_{7}^{2}\right )  | \bx | ^{-2n-2}+\left ( 2\xi_{5}\xi_{6}- 2\xi_{7}\xi_{8}\right )|\bx| ^{-2n} +\left ( \xi_{6}^{2} - \xi_{8}^{2}\right )|\bx|^{-2n+2}\right)\notag\\
		%&\times \omega^{2n-2}\left(1+ \mathcal{O}(\omega^{2})\right)\mathrm{d}\bx\notag\\
		%=&\int_0^{2\pi}\int_{\gamma_2}^{R } \Big\{\left (\left ( \xi_{5}^{2} - \xi_{7}^{2}\right ) r ^{-2n-2}+\left ( 2\xi_{5}\xi_{6}- 2\xi_{7}\xi_{8}\right )r^{-2n} +\left ( \xi_{6}^{2} - \xi_{8}^{2}\right ) r ^{-2n+2}\right )\omega^{2n-2}\notag\\
		%&\times \left(1+\mathcal{O}(\omega^{2})\right)\Big\}r\mathrm{d}r\mathrm{d}\theta\notag\\
		%=& \bigg(\frac{2\pi (\xi_{5}^{2}- \xi_{7}^{2})}{2n} (\gamma_2 ^{-2n}-R^{-2n})+\frac{2\pi (2\xi_{5}\xi_{6}- 2\xi_{7}\xi_{8})}{2n-2} (\gamma_2 ^{-2n+2}-R^{-2n+2})\notag\\
		%&+\frac{2\pi (\xi_{6}^{2}- \xi_{8}^{2})}{2n-4}\left( \gamma_2 ^{-2n+4}-R^{-2n+4}\right)\bigg)\omega^{2n-2}\left(1+ \mathcal{O}(\omega^{2})\right)\notag\\
		=&\left(\zeta _{6,n} (\gamma_2 ^{-2n}-R^{-2n})+\zeta _{7,n} (\gamma_2 ^{-2n+2}-R^{-2n+2})+\zeta _{8,n} \left( \gamma_2 ^{-2n+4}-R^{-2n+4}\right)\right)\notag\\
		& \times  \omega ^{2n-2}\left(1+ \mathcal{O}(\omega^{2})\right)\notag\\
		=&\zeta _{6,n}\frac{R^{2n}-\gamma_2 ^{2n}}{\gamma_2 ^{2n}R^{2n}}\left(1+\zeta _{9,n} \frac{R^{2n-2}-\gamma_{2}^{2n-2}}{R^{2n}-\gamma_2 ^{2n}}\gamma_2 ^{2}R^{2}+\zeta _{10,n} \frac{R^{2n-4}-\gamma_{2}^{2n-4}}{R^{2n}-\gamma_2 ^{2n}}\gamma_2 ^{4}R^{4}\right)\notag\\
		&\times \omega ^{2n-2}\left(1+ \mathcal{O}(\omega^{2})\right),\notag
	\end{align}
	where
	\begin{align}
		&\zeta _{6,n}=\frac{\pi\kappa^2 2^{3-2 n} \left(\tau ^2-1\right)^2 \mu ^{1-n} \rho ^{n-1} \left(\tau ^2 (\lambda +3 \mu )+(n+1) (\lambda +\mu )\right)^2}{n(n+1)^2 \left(\tau ^2+1\right)^2 (\lambda +3 \mu )^2  (n-2)!^2},\notag\\
		&\zeta_{7,n}=-\frac{\pi\kappa^2 4^{2-n} \left(\tau ^2-1\right)^2 (\lambda +\mu ) \mu ^{1-n} \rho ^{n-1} \left(\tau ^2 (\lambda +3 \mu )+(n+1) (\lambda +\mu )\right)}{(n-1)(n+1) \left(\tau ^2+1\right)^2 (\lambda +3 \mu )^2  (n-2)!^2},\notag\\
		&\zeta_{8,n}=\frac{\pi\kappa^2 4^{1-n} \left(\tau ^2-1\right)^2 \rho ^{n-1} \left((\lambda  (n-2)+\mu  (n-4))^2+(2 \mu +n (\lambda +\mu ))^2\right)}{\mu ^{n-1} (n-2)\left(\tau ^2+1\right)^2 (\lambda +3 \mu )^2 (n-1)!^2},\notag\\
		&\zeta _{9,n}=\frac{-2 n (n+1) (\lambda +\mu )}{\left(n^2-1\right) (\lambda +\mu )+(n-1) \tau ^2 (\lambda +3 \mu )}=-2+\frac{\zeta_{\frac{1}{n}}}{n}+\mathcal{O}\left(\frac{1}{n^2}\right),\notag\\
		&\zeta _{10,n}=\frac{n (n+1)^2 \left((\lambda  (n-2)+\mu  (n-4))^2+(2 \mu +n (\lambda +\mu ))^2\right)}{2 (n-2) (n-1)^2 \left(\tau ^2 (\lambda +3 \mu )+(n+1) (\lambda +\mu )\right)^2}\notag\\
		&\quad\quad=1-\frac{\zeta_{\frac{1}{n}}}{n}+\mathcal{O}\left(\frac{1}{n^2}\right),\notag\\
		&\zeta_{\frac{1}{n}}=\frac{2\tau ^2 (\lambda +3 \mu )-2\lambda -2\mu}{(\lambda +\mu )}.\notag
	\end{align}
	Through rigorous calculations, we can derive	
	\begin{align}\label{eq:zeta2}
		\zeta _{6,n}+\zeta _{7,n}+\zeta _{8,n}=\mathcal{O}\left(\frac{4^{1-n} \mu ^{1-n} \rho ^{n-1}}{n  (n-1)!^2}\right),\quad
		1+\zeta _{9,n}+\zeta _{10,n}=\frac{\mathcal{R}_{2}}{n^2}+\mathcal{O}\left(\frac{1}{n^3}\right),
		%+\frac{\mathcal{R}_{3}}{n^3}+\mathcal{O}\left(\frac{1}{n^4}\right),
	\end{align}
	where
	\begin{align}\label{eq:R2}
		\mathcal{R}_{2}=&\frac{3 \lambda ^2+\tau ^4 (\lambda +3 \mu )^2-2 \tau ^2 (\lambda +\mu ) (\lambda +3 \mu )+10 \lambda  \mu +11 \mu ^2}{(\lambda +\mu )^2}.
		%\mathcal{R}_{3}=&\frac{1}{(\lambda +\mu )^4}\big(2 (\lambda +\mu ) (-\tau ^2 (\lambda +3 \mu ) \left(3 \lambda ^2+10 \lambda  \mu +11 \mu ^2\right)\notag\\
		%&+(\lambda +\mu ) \left(5 \lambda ^2+18 \lambda  \mu +21 \mu ^2\right)+\tau ^6 \left(-(\lambda +3 \mu )^3\right)+\tau ^4 (\lambda +\mu ) (\lambda +3 \mu )^2)\big).\notag
	\end{align}
	
	Finally, we show that the external scattered field $\bu^s$ is also concentrated, in terms of energy, on the exterior surface of the domain $D$,
	\begin{align}
		&\frac{\|\mathbf{u}^s\|_{L^2\left(\left(D_R \setminus \overline{D}\right) \setminus \mathcal{N}_{D,\gamma_{2}-1}^{+}\right)^2}^2}{\|\mathbf{u}^s\|_{L^2\left(D_{R}\setminus \overline{D}\right)^2}^2}\notag\\
		%=&\frac{\left(1+\zeta _{9,n} \frac{R^{2n-2}-\gamma_{2}^{2n-2}}{R^{2n}-\gamma_2 ^{2n}}\gamma_2 ^{2}R^{2}+\zeta _{10,n} \frac{R^{2n-4}-\gamma_{2}^{2n-4}}{R^{2n}-\gamma_2 ^{2n}}\gamma_2 ^{4}R^{4}\right)\omega ^{2n-2}\left(1+ \mathcal{O}(\omega^{2})\right)}{\left(1+\zeta _{9,n} \frac{R^{2n-2}-1}{R^{2n}-1}R^{2}+\zeta _{10,n} \frac{R^{2n-4}-1}{R^{2n}-1}R^{4}\right)\omega ^{2n-2}\left(1+ \mathcal{O}(\omega^{2})\right)}\notag\\
		% &\times\frac{\zeta _{6,n}\frac{R^{2n}-\gamma_2 ^{2n}}{\gamma_2 ^{2n}R^{2n}}}{\zeta _{6,n}\frac{R^{2n}-1}{R^{2n}}}\notag\\
		%=&\frac{\left(1-\left(\frac{\gamma_2}{R}\right) ^{2n}\right)\left(1+\left(-2+\mathcal{O}\left(\frac{1}{n}\right)\right)\gamma_{2}^{2}\frac{1-\left(\frac{\gamma_2}{R}\right) ^{2n-2}}{1-\left(\frac{\gamma_{2}}{R}\right) ^{2n}}+\left(1+\mathcal{O}\left(\frac{1}{n}\right)\right)\gamma_{2}^{4}\frac{1-\left(\frac{\gamma_2}{R}\right) ^{2n-4}}{1-\left(\frac{\gamma_{2}}{R}\right) ^{2n}}\right)}{\gamma_{2}^{2n}\left(1-\left(\frac{1}{R}\right) ^{2n}\right)\left(1+\zeta_{9,n}\frac{1-\left(\frac{1}{R}\right)^{2n-2}}{1-\left(\frac{1}{R}\right)^{2n}}+\zeta_{10,n}\frac{1-\left(\frac{1}{R}\right)^{2n-4}}{1-\left(\frac{1}{R}\right)^{2n}}\right)}\notag\\
		%&\times \left(1+\mathcal{O}\left(\omega^2\right)\right)\left(1-\mathcal{O}\left(\omega^2\right)\right)\notag\\
		%=&\frac{1}{\gamma_{2}^{2n}}\left(1-\left(\frac{\gamma_2}{R}\right) ^{2n}\right)\left(1+\frac{1}{R^{2n}}+\mathcal{O}\left(\frac{1}{R^{4n}} \right)\right)\times\frac{1-2\gamma_{2}^2+\gamma_{2}^4+\mathcal{O}\left(\frac{1}{n}\right)}{1+\zeta_{9,n}+\zeta_{10,n}+\mathcal{O}\left(\frac{1}{R^{2n-4}}\right)}\notag\\
		%&\times \left(1+\mathcal{O}\left(\omega^2\right)\right)\left(1-\mathcal{O}\left(\omega^2\right)\right)\notag\\
		%=&\frac{1}{\gamma_{2}^{2n}}\left(1-\left(\frac{\gamma_2}{R}\right) ^{2n}\right)\left(1+\frac{1}{R^{2n}}+\mathcal{O}\left(\frac{1}{R^{4n}} \right)\right)\times\frac{1-2\gamma_{2}^2+\gamma_{2}^4+\mathcal{O}\left(\frac{1}{n}\right)}{\frac{\mathcal{R}_{2}}{n^2}+\mathcal{O}\left(\frac{1}{R^{2n-4}}\right)}\notag\\ 
		%&\times \left(1+\mathcal{O}\left(\omega^2\right)\right)\left(1-\mathcal{O}\left(\omega^2\right)\right)\notag\\
		%=&\frac{n^2}{\mathcal{R}_2\gamma_{2}^{2n}}\left(1-\left(\frac{\gamma_2}{R}\right) ^{2n}\right)\left(1+\frac{1}{R^{2n}}+\mathcal{O}\left(\frac{1}{R^{4n}} \right)\right)\left(1-2\gamma_{2}^2+\gamma_{2}^4+\mathcal{O}\left(\frac{1}{n}\right)\right)\notag\\
		%&\times\left(1-\mathcal{O}\left(\frac{n^2}{R^{2n-3}}\right)\right) \left(1+\mathcal{O}\left(\omega^2\right)\right)\left(1-\mathcal{O}\left(\omega^2\right)\right)\notag\\
		=&\Big\{\frac{1-2\gamma_{2}^2+\gamma_{2}^4}{\mathcal{R}_{2}}\times\frac{n^2}{\gamma_{2}^{2n}}+\mathcal{O}\left(\frac{n}{\gamma_{2}^{2n}}\right)\Big\} \left(1+\mathcal{O}\left(\omega^2\right)\right)\left(1-\mathcal{O}\left(\omega^2\right)\right)\notag\\
		\leq&K_{2}\times\frac{n^2}{\gamma_{2}^{2n}} \left(1+\mathcal{O}\left(\omega^2\right)\right)\left(1-\mathcal{O}\left(\omega^2\right)\right)\notag\\
		\leq&\gamma_2^{-n}\left(1+\mathcal{O}\left(\omega^2\right)\right)\left(1-\mathcal{O}\left(\omega^2\right)\right)\leq\varepsilon+\mathcal{O}\left(\varepsilon\omega^2\right)\ll1,\notag
	\end{align}
	where $K_{2} =\frac{2\left(1-2\gamma_{2}^2+\gamma_{2}^4\right)}{\mathcal{R}_{2}}$ is defined in \eqref{eq:lemk1k2}. According to Lemma \ref{lem:thmxy}, the inequality $K_{2}n^2 \leq \gamma_{2}^{n}$ holds.
	
	The proof is complete.
\end{proof}

\end{document}
